\section{Fragmentos de código de implementación} \label{sec:anexo_fragmentos_codigo}

Este anexo reúne fragmentos representativos de la implementación utilizada durante el desarrollo de la metodología descrita en el Capítulo~\ref{sec:obtencion_datos_section}. En el cuerpo principal de la tesis se privilegia el uso de pseudocódigo y descripciones de alto nivel; en cambio, aquí se presentan extractos concretos del código para facilitar trazabilidad e interpretación técnica.\\

La organización de este anexo sigue la misma estructura del capítulo de Metodología: obtención de datos, benchmark experimental, construcción del grafo y modelos basados en GNN.\\

\subsection{Obtención de Datos}

\subsubsection{Extracción de rutas BGP con BGPStream}

\begin{lstlisting}[language=Python, caption=Implementación de descarga y filtrado de rutas BGP con BGPStream, label=anxcode:bgpstream_download]
import datetime
import pybgpstream

def downloader(start_date, duration):
    base = int(datetime.datetime.strptime(start_date, "%m/%d/%Y").strftime("%s"))

    stream = pybgpstream.BGPStream()
    stream.add_interval_filter(base, base + int(duration))
    stream.add_filter("record-type", "ribs")
    stream.start()

    path_set = set()
    f = open("data/rib.txt", "w")

    rec = stream.get_next_record()
    while rec:
        elem = rec.get_next_elem()
        while elem:
            path = elem.fields["as-path"]
            if "{" in path or "(" in path:
                elem = rec.get_next_elem()
                continue

            prefix = elem.fields["prefix"]
            if ":" not in prefix and path not in path_set:
                f.write(path.replace(" ", "|") + "\\n")
                path_set.add(path)

            elem = rec.get_next_elem()

        rec = stream.get_next_record()

    f.close()
\end{lstlisting}

\subsubsection{Carga estandarizada de rutas BGP}

\begin{lstlisting}[language=Python, caption=Implementación de importación de rutas BGP desde archivo, label=anxcode:importar_bgp_routes]
bgp_routes = [] 

with open(nombre_archivo, 'r') as archivo:
    for i,linea in enumerate(archivo):
        if i > MAX_NUM_ROUTES:
            break
        numeros = linea.strip().split('|')
        bgp_routes.append(numeros) 

print("[RUTAS BGP]:",bgp_routes[:10])
print("[CANTIDAD PATHS]",len(bgp_routes))
\end{lstlisting}

\subsection{Benchmark Experimental}

\subsubsection{Transformación y limpieza de predicciones}

\begin{lstlisting}[language=Python, caption=Conversión de etiquetas de Gao al formato estándar, label=anxcode:transformacion_estandar_gao]
y_test_prediction_new = []
for i in range(len(y_test_prediction)):
    if y_test_prediction[i] % 2 == 0:
        y_test_prediction_new.append(0)
    elif y_test_prediction[i] == 3:
        y_test_prediction_new.append(2)
    elif y_test_prediction[i] == 1:
        y_test_prediction_new.append(1)
    else:
        y_test_prediction_new.append(-1)

y_test_prediction_new = np.asarray(y_test_prediction_new)
print(len(y_test_prediction_new))
y_test_prediction = y_test_prediction_new
\end{lstlisting}

\begin{lstlisting}[language=Python, caption=Limpieza de relaciones sin ground truth válido, label=anxcode:limpieza_valores_no_encontrados]
x_test_cleaned = np.asarray([np.asarray(x_test[i]) for i in range(len(x_test)) if y_test_prediction[i] != -1])
y_test_cleaned = np.asarray([y_test[i] for i in range(len(y_test)) if y_test_prediction[i] != -1])

y_test_prediction_cleaned = np.asarray([y_test_prediction[i] for i in range(len(y_test_prediction)) if y_test_prediction[i] != -1])
\end{lstlisting}

\subsubsection{Integración de BGP2Vec}

\begin{lstlisting}[language=Python, caption=Inicialización del modelo BGP2Vec, label=anxcode:bgp2vec_init]
bgp2vec = BGP2VEC(model_path = MODELS_PATH,
                  bgp_routes = bgp_routes,
                  rewrite = True,
                  test_limit = test_limit,
                  mode = mode,
                  epochs = epochs)
\end{lstlisting}

\begin{lstlisting}[language=Python, caption=Mapeo de ASN a índices para entrenamiento y prueba en BGP2Vec, label=anxcode:bgp2vec_mapping]
x_training_mapped = []
y_training_filtered = []

for (a, b), y in zip(x_training, y_training):
    a_index = asn_to_index.get(str(a), 0)
    b_index = asn_to_index.get(str(b), 0)
    if a_index == 0 or b_index == 0:
        continue
    x_training_mapped.append([a_index, b_index])
    y_training_filtered.append(y)

x_training_mapped = np.array(x_training_mapped)
y_training_filtered = np.array(y_training_filtered)

x_test_mapped = []
y_test_filtered = []

for (a, b), y in zip(x_test, y_test):
    a_index = asn_to_index.get(str(a), 0)
    b_index = asn_to_index.get(str(b), 0)
    if a_index == 0 or b_index == 0:
        continue
    x_test_mapped.append([a_index, b_index])
    y_test_filtered.append(y)

x_test_mapped = np.array(x_test_mapped)
y_test_filtered = np.array(y_test_filtered)
\end{lstlisting}

\subsection{Construcción del Grafo de Internet}

\subsubsection{Construcción del grafo en formato DGL}

\begin{lstlisting}[language=Python, caption=Clase base para construir el dataset del grafo, label=anxcode:graph_class]
class Graph:
    def __init__(self,dataset_graph_path ,max_paths,debug = False):
        self.debug = debug
        self.nx_graph = None
        self.data_path = dataset_graph_path
        
    def create_topology_from_ribs(self, rib_filename:str,filename_out="edges.csv"):
        ...

    def features_nodes(self,features_filename,filename_out="nodes.csv"):
        ...

    def remove_nodes_degree(self, degree,filename_out="edges.csv"):
        ...
\end{lstlisting}

\subsection{Modelos Basados en GNN}

\subsubsection{Modelos para inferencia con GNN y CNN}

\begin{lstlisting}[language=Python, caption=Implementación de un modelo GCN con distintos decodificadores, label=anxcode:modelo_gcn]
class GCN(nn.Module):
    def __init__(self, in_feats, hidden_feats, out_feats, out_feats_mlp=1, drop=0.3):
        super().__init__()
        self.conv1 = GraphConv(in_feats, hidden_feats)
        self.conv2 = GraphConv(hidden_feats, out_feats)
    
        self.MLP = MLPPredictor(out_feats,out_feats_mlp)
        self.BilinearDecoder = BilinearPredictor(out_feats, 3)
        self.regressor = nn.Linear(out_feats, 1)
        self.drop  = nn.Dropout(drop)

    def encode(self, g, x):
        g = dgl.add_self_loop(g)
        h = self.drop(F.relu(self.conv1(g, x)))
        h = self.conv2(g, h)
        return h

    def decodeDotProduct(self, g, h):
        with g.local_scope():
            g.ndata["h"] = h
            g.apply_edges(fn.u_dot_v("h", "h", "score"))
            return g.edata["score"][:, 0]

    def decodeMLP(self, g, h):
        return self.MLP(g, h)

    def decodeBilinear(self, g, h):
        return self.BilinearDecoder(g, h)

    def forward(self, g, x):
        h = self.encode(g, x)
        out = self.regressor(h).squeeze(-1)
        return out
\end{lstlisting}

\begin{lstlisting}[language=Python, basicstyle=\ttfamily\footnotesize, caption=Mapeo final de ASN a índices para la etapa supervisada, label=anxcode:mapeo_asn_id]
x_training_mapped = []
y_training_filtered = []

for (a, b), y in zip(x_training, y_training):
    a_index = asn_to_index.get(str(a), 0)
    b_index = asn_to_index.get(str(b), 0)

    if a_index == 0 or b_index == 0:
        continue

    x_training_mapped.append([a_index, b_index])
    y_training_filtered.append(y)

x_training_mapped = np.array(x_training_mapped)
y_training_filtered = np.array(y_training_filtered)

x_test_mapped = []
y_test_filtered = []

for (a, b), y in zip(x_test, y_test):
    a_index = asn_to_index.get(str(a), 0)
    b_index = asn_to_index.get(str(b), 0)

    if a_index == 0 or b_index == 0:
        continue

    x_test_mapped.append([a_index, b_index])
    y_test_filtered.append(y)

x_test_mapped = np.array(x_test_mapped)
y_test_filtered = np.array(y_test_filtered)

np.save(PATH + DATA + "x_training_mapped.npy", x_training_mapped)
np.save(PATH + DATA + "y_training_filtered.npy", y_training_filtered)
np.save(PATH + DATA + "x_test_mapped.npy", x_test_mapped)
np.save(PATH + DATA + "y_test_filtered.npy", y_test_filtered)
\end{lstlisting}

\begin{lstlisting}[basicstyle=\ttfamily, caption=Modelo CNN utilizado en el enfoque BGP2Vec, label=anxcode:cnn_bgp2vec_model]
experiment = None
num_classes = len(class_names)
embedding_trainable = False
input_length = 2
embedding_vector_length = embeddings.shape[1]
total_ASNs = embeddings.shape[0]

model = Sequential()
model.add(Embedding(total_ASNs, embedding_vector_length, input_length=input_length,
                    weights=[embeddings], trainable=embedding_trainable))
model.add(Conv1D(filters=32, kernel_size=3, padding='same', activation='relu'))
model.add(Reshape((32, 2)))
model.add(MaxPooling1D(pool_size=2))
model.add(Conv1D(filters=32, kernel_size=3, padding='same', activation='relu'))
model.add(Reshape((32, 16)))
model.add(MaxPooling1D(pool_size=2))
model.add(Flatten())
model.add(Dense(100, activation='relu'))
model.add(Dense(num_classes, activation='softmax'))
\end{lstlisting}

\begin{lstlisting}[basicstyle=\ttfamily, caption=Cálculo de pesos por clase para balanceo, label=anxcode:class_weights_dict]
from sklearn.utils import class_weight
import numpy as np

class_weights = class_weight.compute_class_weight(
                    class_weight='balanced', 
                    classes=np.array(list(range(num_classes))), 
                    y=y_training_filtered)

class_weight_dict = {i: w for i, w in enumerate(class_weights)}
print('[class_weights_dict]:', class_weight_dict)
\end{lstlisting}